\documentclass[conference]{IEEEtran}
\IEEEoverridecommandlockouts
\usepackage{cite}
\usepackage{amsmath,amssymb,amsfonts}
\usepackage{algorithmic}
\usepackage{graphicx}
\usepackage{textcomp}
\usepackage{xcolor}
\usepackage{listings}

\begin{document}

\title{Low-Level Systems Programming in Terminal Game Development: A Custom Memory and Rendering Approach in C}

\author{\IEEEauthorblockN{Ayush Gupta}
\IEEEauthorblockA{\textit{Department of Computer Science} \\
\textit{University Name}\\
City, Country \\
email@example.com}
}

\maketitle

\begin{abstract}
This paper discusses the implementation of a terminal-based Snake game developed in C, focusing on low-level systems programming techniques. Unlike conventional implementations that rely on high-level libraries like ncurses or the standard C library for memory management, this project implements a custom linked-list based memory allocator, raw ANSI escape sequence rendering, and non-blocking I/O handling via termios.
\end{abstract}

\section{Introduction}
Modern software development often abstracts hardware and OS interactions through extensive libraries. This project aims to strip back these layers by building a classic Snake game using only core POSIX APIs and custom logic.

\section{System Architecture}
\subsection{Custom Memory Management}
The core of the memory system is a static character array of 1MB. A custom allocator (\textit{alloc}) and deallocator (\textit{dealloc}) manage this space using a free-list approach. 

\subsection{Terminal Rendering Engine}
Rendering is achieved through ANSI escape sequences, allowing for high-speed updates without the overhead of UI libraries.

\subsection{Non-Blocking Input Handling}
The terminal settings are modified via \texttt{termios} to capture immediate keypresses, enabling real-time gameplay.

\section{Conclusion}
The project successfully demonstrates complex systems programming concepts within a compact, performant application.

\end{document}
